\chapter{Programs — Emod}
%%% generated by the blueprint pipeline from TCSlib.Programs.Light.Lib.Emod
\section{Overview}

Declarations of the Lean module \texttt{TCSlib.Programs.Light.Lib.Emod}.

\section{Declarations}

\begin{definition}[Argument index]
\lean{Light.Emod.Arg}
\label{Light.Emod.Arg}
\leanok
The argument index is the natural number $0$.
\end{definition}

\begin{definition}[Modulus index]
\lean{Light.Emod.Modulus}
\label{Light.Emod.Modulus}
\leanok
\uses{Light.Emod.Arg}
The modulus index is the natural number $1$.
\end{definition}

\begin{definition}[Free-variable index]
\lean{Light.Emod.Free}
\label{Light.Emod.Free}
\leanok
\uses{Light.Emod.Arg}
The free-variable index is the natural number $2$.
\end{definition}

\begin{definition}[Rest-variable index]
\lean{Light.Emod.Rest}
\label{Light.Emod.Rest}
\leanok
\uses{Light.Emod.Arg}
The rest-variable index is the natural number $3$.
\end{definition}

\begin{definition}[Count-variable index]
\lean{Light.Emod.Count}
\label{Light.Emod.Count}
\leanok
\uses{Light.Emod.Arg}
The count-variable index is the natural number $4$.
\end{definition}

\begin{definition}[Multiplication-variable index]
\lean{Light.Emod.Mult}
\label{Light.Emod.Mult}
\leanok
\uses{Light.Emod.Arg}
The multiplication-variable index is the natural number $5$.
\end{definition}

\begin{definition}[Emod table round]
\lean{Light.emodTableRound}
\label{Light.emodTableRound}
\leanok
\uses{Light.Emod.Count, Light.Emod.Free, Light.Emod.Mult, Light.Stmt, Light.k, Light.v}
A table round stores the product of `Mult` and the sum of `Free` and `Count` at the
memory address given by that sum, then increments `Count` by one and doubles `Mult`.
\end{definition}

\begin{definition}[Emod table loop]
\lean{Light.emodTable}
\label{Light.emodTable}
\leanok
\uses{Light.Emod.Mult, Light.Emod.Rest, Light.Stmt, Light.emodTableRound, Light.v}
The table loop repeats the table round while `Mult` is at most `Rest`.
\end{definition}

\begin{definition}[Emod reduction round]
\lean{Light.emodReduceRound}
\label{Light.emodReduceRound}
\leanok
\uses{Light.Emod.Count, Light.Emod.Free, Light.Emod.Rest, Light.M, Light.Stmt, Light.k, Light.v}
A reduction round decrements `Count` by one, using truncated subtraction on natural
numbers. It then subtracts the product of `M` and the sum of `Free` and the updated
`Count` from `Rest` if that product is at most `Rest`; otherwise it leaves `Rest`
unchanged.
\end{definition}

\begin{definition}[Emod reduction loop]
\lean{Light.emodReduce}
\label{Light.emodReduce}
\leanok
\uses{Light.Emod.Count, Light.Stmt, Light.emodReduceRound, Light.k, Light.v}
The reduction loop repeats a reduction round while `Count` is positive.
\end{definition}

\begin{definition}[Emod sign adjustment]
\lean{Light.emodSign}
\label{Light.emodSign}
\leanok
\uses{Light.Emod.Arg, Light.Emod.Modulus, Light.Emod.Rest, Light.Stmt, Light.k, Light.v}
If `Arg` is negative, set it to zero when `Rest` is zero, and otherwise set it to
`Modulus` minus `Rest`, using truncated subtraction on natural numbers. If `Arg` is
nonnegative, set it to `Rest`.
\end{definition}

\begin{definition}[Emod algorithm body]
\lean{Light.emodBody}
\label{Light.emodBody}
\leanok
\uses{Light.Emod.Arg, Light.Emod.Count, Light.Emod.Modulus, Light.Emod.Mult, Light.Emod.Rest, Light.Stmt, Light.emodReduce, Light.emodSign, Light.emodTable, Light.k, Light.v}
The algorithm sets `Rest` to the absolute value of `Arg`, using truncated subtraction
from zero when `Arg` is negative; it sets `Count` to zero and `Mult` to `Modulus`, then
runs the table loop, the reduction loop, and the sign adjustment in sequence.
\end{definition}

\begin{definition}[Emod round count]
\lean{Light.emodRounds}
\label{Light.emodRounds}
\leanok
For natural numbers $a$ and $M$, the number of rounds is the natural-number size of the
quotient of $a$ by $M$.
\end{definition}

\begin{definition}[Emod running time]
\lean{Light.emodTime}
\label{Light.emodTime}
\leanok
\uses{Light.emodRounds}
For natural numbers $a$ and $M$, the running time is $43$ times the number of rounds,
plus $34$.
\end{definition}

\begin{theorem}[Power threshold and round count]
\lean{Light.mul_two_pow_le_iff}
\label{Light.mul_two_pow_le_iff}
\leanok
\uses{Light.emodRounds}
\difficulty{3}
For natural numbers $a,M,i$ with $1\le M$, the inequality $M2^i\le a$ holds if and only
if $i$ is less than the number of rounds, defined as the binary size of the quotient
$a/M$.
\end{theorem}

\begin{theorem}[Bound on remainder rounds]
\lean{Light.emodRounds_le}
\label{Light.emodRounds_le}
\leanok
\uses{Light.emodRounds}
\difficulty{4}
For natural numbers $a,M,V$ with $a\le V$, the number of base-size rounds needed to
reduce $a$ by division by $M$, defined as the binary size of $\lfloor a/M\rfloor$, is at
most $\lfloor\log_2(V+1)\rfloor+1$.
\end{theorem}

\begin{theorem}[Negated remainder modulo a natural number]
\lean{Light.neg_emod_nat}
\label{Light.neg_emod_nat}
\leanok
\difficulty{5}
For natural numbers $a,M$ with $1\le M$, the remainder of $-a$ modulo $M$, computed in
the integers, is $0$ when the natural remainder of $a$ modulo $M$ is zero; otherwise it
is $M$ minus that remainder.
\end{theorem}

\begin{definition}[Preliminary limit conditions]
\lean{Light.Emod.Pre}
\label{Light.Emod.Pre}
\leanok
\uses{Light.Limits}
Let the word limit be an integer and the space limit, depth limit, and parameters
$M_o,V,f_r$ be natural numbers. The limits and parameters satisfy the preliminary
conditions if $1\le M_o$, the space limit is at most the word limit, $4V+4M_o+16$ is at
most the word limit, and $f_r+\lfloor\log_2(V+1)\rfloor+2$ is at most the space limit.
\end{definition}

\begin{definition}[Table condition]
\lean{Light.Emod.Table}
\label{Light.Emod.Table}
\leanok
\uses{Light.Kept}
Let $\mu,\mu':\mathbb N\to\mathbb Z$ and let $M_o,f_r,i$ be natural numbers. The table
condition holds when, for every natural number $j<i$, $\mu'(f_r+j)=M_o2^j$, and the pair
$(\mu,\mu')$ satisfies the relation called $\mathrm{Kept}$ at offset $f_r$.
\end{definition}

\begin{theorem}[Bounded construction of a doubling table]
\lean{Light.Emod.table_ends}
\label{Light.Emod.table_ends}
\leanok
\uses{Light.Emod.Pre, Light.Emod.Table, Light.Ends, Light.Ends.whileConst, Light.Stmt.blockCost, Light.emodRounds, Light.emodRounds_le, Light.emodTable, Light.emodTableRound, Light.frame, Light.mul_two_pow_le_iff}
\difficulty{6}
Let $\mathrm{lim}$ be a set of machine limits, $P$ a program, $d$ a call depth, $\mu$ a
memory function, and $x$ an integer. Let $M,V,fr,a$ be natural numbers such that $a\le
V$, $1\le M$, the space limit is at most the word limit, $4V+4M+16$ is at most the word
limit, and $fr+\bigl(\lfloor\log_2(V+1)\rfloor+2\bigr)$ is at most the space limit. Set
$r$ equal to the number of binary digits of $\lfloor a/M\rfloor$. Starting with locals
$[x,M,fr,a,0,M]$ and memory $\mu$, the table-building program has an execution of cost
at most $19r+6$ whose final locals are $[x,M,fr,a,r,M2^r]$, and whose final memory
$\mu'$ satisfies $\mu'(fr+j)=M2^j$ for every $j<r$ and agrees with $\mu$ at every
address below $fr$.
\end{theorem}

\begin{theorem}[Bounded execution of remainder reduction]
\lean{Light.Emod.reduce_ends}
\label{Light.Emod.reduce_ends}
\leanok
\uses{Light.Emod.Pre, Light.Emod.Table, Light.Ends, Light.Ends.iteLast, Light.Ends.setTo, Light.Ends.skip, Light.Ends.whileConst, Light.Stmt.blockCost, Light.emodReduce, Light.emodReduceRound, Light.emodRounds, Light.emodRounds_le, Light.frame, Light.mul_two_pow_le_iff}
\difficulty{6}
Let the machine limits be given, and let $P$ be a program and $d$ a call depth. For
natural numbers $M,V,fr,a$ with $a\le V$, assume $1\le M$, the space limit is at most
the word limit, $4V+4M+16$ is at most the word limit, and
$fr+\bigl(\operatorname{Nat.log}(2,V+1)+2\bigr)$ is at most the space limit. Let $r$ be
the binary size of the quotient of $a$ by $M$, and suppose a memory function $\mu'$
stores $M2^j$ at address $fr+j$ for every $j<r$ and agrees with a memory function $\mu$
at every address below $fr$. From the state whose locals are $[x,M,fr,a,r,t]$ and whose
memory is $\mu'$, the remainder-reduction program has an execution of cost at most
$24r+4$ ending in the state whose locals are $[x,M,fr,a\bmod M,0,t]$ and whose memory is
still $\mu'$.
\end{theorem}

\begin{theorem}[Bounded execution of the signed remainder step]
\lean{Light.Emod.sign_ends}
\label{Light.Emod.sign_ends}
\leanok
\uses{Light.Emod.Pre, Light.Ends, Light.Ends.iteLast, Light.Ends.setTo, Light.emodSign, Light.frame, Light.neg_emod_nat}
\difficulty{5}
Let the machine limits be given, along with a program, a call depth, and natural numbers
$M,V,fr,a$. Suppose $1\le M$, the space limit is at most the word limit, $4V+4M+16$ is
at most the word limit, $fr+\bigl(\lfloor\log_2(V+1)\rfloor+2\bigr)$ is at most the
space limit, and $a=|x|$ for an integer $x$. From a state whose locals are initialized
by the frame $[x,M,fr,a\bmod M,0,t]$ and whose memory is any function $\mu'$, the signed
remainder program has an execution of cost at most $12$ ending in a state whose local at
position $0$ is the integer remainder of $x$ modulo $M$ and whose memory remains $\mu'$.
\end{theorem}

\begin{theorem}[Bounded modular reduction]
\lean{Light.emod_meets}
\label{Light.emod_meets}
\leanok
\uses{Light.Emod.Count, Light.Emod.Modulus, Light.Emod.Mult, Light.Emod.Pre, Light.Emod.reduce_ends, Light.Emod.sign_ends, Light.Emod.table_ends, Light.Ends, Light.Ends.iteThen, Light.Ends.mono, Light.Ends.next, Light.Ends.setToThen, Light.Kept, Light.Meets, Light.Meets.of_body, Light.emodBody, Light.emodReduce, Light.emodRounds, Light.emodSign, Light.emodTable, Light.emodTime, Light.frame, Light.k, Light.v}
\difficulty{5}
Let $P$ be a program with word, space, and depth limits, and let $p$ be a procedure
index. Suppose procedure $p$ is the modular-reduction body. Let $M,V,fr$ be natural
numbers such that $1\le M$, the space limit is at most the word limit, $4V+4M+16$ is at
most the word limit, and $fr+\bigl(\lfloor\log_2(V+1)\rfloor+2\bigr)$ is at most the
space limit. For any initial memory $\mu$ and integer $x$ with $|x|\le V$, running this
procedure at depth $d$ on the arguments $(x,M,fr)$ terminates within $43$ times the
binary bit length of $\lfloor |x|/M\rfloor$, plus $34$, steps. It returns the remainder
of $x$ modulo $M$ and leaves every memory location below $fr$ unchanged.
\end{theorem}
